\begin{proof}[Proof of \cref{obj:lem:isotropic-row-signing}]
For a deterministic row tuple \(\boldsymbol v=(v_1,\ldots,v_n)\in(\mathbb R^{r_0})^n\), define
\[
\begin{aligned}
\mathsf D_{\rm row}(\boldsymbol v)
&:=\sum_{i=1}^n\|v_i\|_2^2,\\
\mathsf O_{\rm row}(\boldsymbol v,z)
&:=\sum_{\substack{1\leq i,i'\leq n\\i\ne i'}}
z_i z_{i'}\langle v_i,v_{i'}\rangle,
\qquad z\in\{-1,1\}^n,\\
\mathsf M_{\rm off}(\boldsymbol v)
&:=\max_{z\in\{-1,1\}^n}
|\mathsf O_{\rm row}(\boldsymbol v,z)|,\\
\mathsf M_{\rm sign}(\boldsymbol v)
&:=\min_{z\in\{-1,1\}^n}
\left\|\sum_{i=1}^n z_i v_i\right\|_2^2.
\end{aligned}
\]
For \(n=0\), the sign space consists of the empty tuple, and all four quantities are zero. Square-integrability ensures integrability of the finite extrema used below.

\begin{enumerate}
\item Suppose first that \(n\leq256\).
Expanding the squared norm gives, for every signing,
\[
\left\|\sum_{i=1}^n z_i v_i\right\|_2^2
=\mathsf D_{\rm row}(\boldsymbol v)
+\mathsf O_{\rm row}(\boldsymbol v,z).
\]
Consequently,
\[
\mathsf M_{\rm sign}(\boldsymbol v)
\geq\mathsf D_{\rm row}(\boldsymbol v)
-\mathsf M_{\rm off}(\boldsymbol v).
\]
Isotropy gives
\[
\mathbb E\|v\|_2^2
=\sum_{\alpha=1}^{r_0}\mathbb E(v_\alpha^2)=r_0,
\qquad
\mathbb E\mathsf D_{\rm row}(\boldsymbol v)=nr_0,
\]
and hence
\[
\mathbb E\mathsf M_{\rm sign}(\boldsymbol v)
\geq nr_0-\mathbb E\mathsf M_{\rm off}(\boldsymbol v).
\]
These identities and the diagonal-subtraction inequality hold for every \(n\).
% lean: CausalSmith.Experimentation.BivariateSobolevDesignCapacity.isotropic_signing_lower_of_offDiagonal

For distinct rows, independence and isotropy yield
\[
\begin{aligned}
\mathbb E\langle v_i,v_{i'}\rangle^2
&=\sum_{\alpha,\beta=1}^{r_0}
\mathbb E[(v_i)_\alpha(v_i)_\beta]\,
\mathbb E[(v_{i'})_\alpha(v_{i'})_\beta]\\
&=\sum_{\alpha,\beta=1}^{r_0}\mathbf 1_{\{\alpha=\beta\}}
=r_0.
\end{aligned}
\]
Thus Cauchy--Schwarz implies
\[
\mathbb E|\langle v_i,v_{i'}\rangle|
\leq\sqrt{r_0}.
\]
The triangle inequality, followed by the loose count of at most \(n^2\) ordered pairs, gives
\[
\mathsf M_{\rm off}(\boldsymbol v)
\leq\sum_{\substack{1\leq i,i'\leq n\\i\ne i'}}
|\langle v_i,v_{i'}\rangle|,
\qquad
\mathbb E\mathsf M_{\rm off}(\boldsymbol v)
\leq n^2\sqrt{r_0}.
\]
Therefore,
\[
\mathbb E\mathsf M_{\rm sign}(\boldsymbol v)
\geq nr_0-n^2\sqrt{r_0}.
\]
% lean: CausalSmith.Experimentation.BivariateSobolevDesignCapacity.isotropic_signing_lower_pair_count

Since \(n\leq256\),
\[
\sqrt n\leq16,\qquad
n=(\sqrt n)^2\leq16\sqrt n,
\qquad
n^2\sqrt{r_0}\leq16n\sqrt{nr_0}.
\]
It follows that
\[
\mathbb E\mathsf M_{\rm sign}(\boldsymbol v)
\geq nr_0-16n\sqrt{nr_0}
\geq nr_0-n^2-16n\sqrt{nr_0}.
\]
The same calculation covers \(n=0\).
% lean: CausalSmith.Experimentation.BivariateSobolevDesignCapacity.isotropic_small_sample_signing_lower

\item Suppose now that \(n>256\). We first establish the finite-class estimate needed for the partition argument.
Let
\[
n_{\rm loc}\in\mathbb N_0,\qquad
\boldsymbol v_{\rm loc}
=(v^{\rm loc}_1,\ldots,v^{\rm loc}_{n_{\rm loc}})
\in(\mathbb R^{r_0})^{n_{\rm loc}},
\]
where the local rows are independent copies of \(v\). Let
\[
\varnothing\ne\mathcal U_{\rm fin}\subset\mathbb R^{r_0}
\quad\text{be finite},\qquad
R_{\rm dir}:=\max_{u_{\rm dir}\in\mathcal U_{\rm fin}}
\|u_{\rm dir}\|_2.
\]
For \(u_{\rm dir}\in\mathcal U_{\rm fin}\), define
\[
\mathsf f_{u_{\rm dir}}:\mathbb R^{r_0}\to\mathbb R,
\qquad
\mathsf f_{u_{\rm dir}}(v):=|\langle u_{\rm dir},v\rangle|,
\]
and
\[
\mathsf C_{\rm fin}(\boldsymbol v_{\rm loc})
:=
\max_{u_{\rm dir}\in\mathcal U_{\rm fin}}
\left|
\sum_{i=1}^{n_{\rm loc}}
\bigl(\mathsf f_{u_{\rm dir}}(v_i^{\rm loc})
-\mathbb E\mathsf f_{u_{\rm dir}}(v)\bigr)
\right|.
\]
Isotropy and Cauchy--Schwarz give
\[
\mathbb E\langle u_{\rm dir},v\rangle^2
=u_{\rm dir}^{\mathsf T}\mathbb E[vv^{\mathsf T}]u_{\rm dir}
=\|u_{\rm dir}\|_2^2,
\qquad
\mathbb E\mathsf f_{u_{\rm dir}}(v)
\leq\|u_{\rm dir}\|_2.
\]
Adding and subtracting the means therefore yields the pointwise bound
\[
\max_{u_{\rm dir}\in\mathcal U_{\rm fin}}
\sum_{i=1}^{n_{\rm loc}}
\mathsf f_{u_{\rm dir}}(v_i^{\rm loc})
\leq n_{\rm loc}R_{\rm dir}
+\mathsf C_{\rm fin}(\boldsymbol v_{\rm loc}).
\]

Introduce an independent copy
\[
\boldsymbol v_{\rm ghost}
=(v^{\rm ghost}_1,\ldots,v^{\rm ghost}_{n_{\rm loc}})
\]
of the local sample, and independent fair signs
\[
\rho=(\rho_1,\ldots,\rho_{n_{\rm loc}})
\in\{-1,1\}^{n_{\rm loc}},
\qquad
\mathbb P(\rho_i=1)=\mathbb P(\rho_i=-1)=\tfrac12.
\]
Define
\[
\mathsf S_\rho(\boldsymbol v_{\rm loc})
:=\sum_{i=1}^{n_{\rm loc}}\rho_i v_i^{\rm loc}.
\]
Jensen's inequality, exchanging the two entries of each independent sample pair according to \(\rho_i\), and the triangle inequality give
\[
\begin{aligned}
\mathbb E\mathsf C_{\rm fin}(\boldsymbol v_{\rm loc})
&\leq
\mathbb E\max_{u_{\rm dir}\in\mathcal U_{\rm fin}}
\left|\sum_{i=1}^{n_{\rm loc}}
\bigl(\mathsf f_{u_{\rm dir}}(v_i^{\rm loc})
-\mathsf f_{u_{\rm dir}}(v_i^{\rm ghost})\bigr)\right|\\
&=
\mathbb E\max_{u_{\rm dir}\in\mathcal U_{\rm fin}}
\left|\sum_{i=1}^{n_{\rm loc}}\rho_i
\bigl(\mathsf f_{u_{\rm dir}}(v_i^{\rm loc})
-\mathsf f_{u_{\rm dir}}(v_i^{\rm ghost})\bigr)\right|\\
&\leq
2\mathbb E\max_{u_{\rm dir}\in\mathcal U_{\rm fin}}
\left|\sum_{i=1}^{n_{\rm loc}}\rho_i
\mathsf f_{u_{\rm dir}}(v_i^{\rm loc})\right|.
\end{aligned}
\]
% lean: CausalSmith.Experimentation.BivariateSobolevDesignCapacity.finite_first_moment_symmetrization

Apply the contraction inequality of
\citep[Theorem~12(4), with Rademacher complexity as in Definition~2]{BartlettMendelson2002Complexities}
to the finite class of linear projections and the map
\[
\varphi_{\rm abs}:\mathbb R\to\mathbb R,
\qquad
\varphi_{\rm abs}(t_{\rm abs}):=|t_{\rm abs}|.
\]
This map is \(1\)-Lipschitz and vanishes at zero. The class is its own countable dense subclass, and a common integrable envelope is
\[
\mathsf E_{\rm fin}(v)
:=\sum_{u_{\rm dir}\in\mathcal U_{\rm fin}}
|\langle u_{\rm dir},v\rangle|
\leq|\mathcal U_{\rm fin}|R_{\rm dir}\|v\|_2,
\qquad
\mathbb E\|v\|_2\leq\sqrt{r_0}.
\]
The envelope bound holds at every \(v\). Contraction supplies
\[
\mathbb E\max_{u_{\rm dir}\in\mathcal U_{\rm fin}}
\left|\sum_{i=1}^{n_{\rm loc}}\rho_i
|\langle u_{\rm dir},v_i^{\rm loc}\rangle|\right|
\leq
2\mathbb E\max_{u_{\rm dir}\in\mathcal U_{\rm fin}}
|\langle \mathsf S_\rho(\boldsymbol v_{\rm loc}),u_{\rm dir}\rangle|.
\]
Combining the preceding bounds proves
\[
\mathbb E\max_{u_{\rm dir}\in\mathcal U_{\rm fin}}
\sum_{i=1}^{n_{\rm loc}}|\langle u_{\rm dir},v_i^{\rm loc}\rangle|
\leq n_{\rm loc}R_{\rm dir}
+4\mathbb E\max_{u_{\rm dir}\in\mathcal U_{\rm fin}}
|\langle \mathsf S_\rho(\boldsymbol v_{\rm loc}),u_{\rm dir}\rangle|.
\]
All sums vanish when \(n_{\rm loc}=0\).
% lean: CausalSmith.Experimentation.BivariateSobolevDesignCapacity.finite_projection_sum_integral_le_exact_rad

\item Consider two independent groups of rows,
\[
\begin{gathered}
n_{\rm left},n_{\rm right}\in\mathbb N_0,\\
\boldsymbol v_{\rm left}
=(v^{\rm left}_1,\ldots,v^{\rm left}_{n_{\rm left}}),
\qquad
\boldsymbol v_{\rm right}
=(v^{\rm right}_1,\ldots,v^{\rm right}_{n_{\rm right}}),
\end{gathered}
\]
each consisting of independent copies of \(v\). For any local row tuple, define
\[
\mathsf K_{\rm row}(\boldsymbol v_{\rm loc})
:=
\max_{z_{\rm loc}\in\{-1,1\}^{n_{\rm loc}}}
\left\|\sum_{i=1}^{n_{\rm loc}}
(z_{\rm loc})_i v_i^{\rm loc}\right\|_2.
\]
Also define the finite direction set and cross-group quantity
\[
\begin{aligned}
\mathcal U_{\rm right}(\boldsymbol v_{\rm right})
&:=
\left\{\sum_{i'=1}^{n_{\rm right}}
(z_{\rm right})_{i'}v_{i'}^{\rm right}:
z_{\rm right}\in\{-1,1\}^{n_{\rm right}}\right\},\\
\mathsf H_{\rm cross}(\boldsymbol v_{\rm left},\boldsymbol v_{\rm right})
&:=
\max_{u_{\rm dir}\in\mathcal U_{\rm right}(\boldsymbol v_{\rm right})}
\sum_{i=1}^{n_{\rm left}}
|\langle u_{\rm dir},v_i^{\rm left}\rangle|.
\end{aligned}
\]
For an empty tuple, \(\mathsf K_{\rm row}=0\); if the right group is empty, \(\mathcal U_{\rm right}=\{0\}\). Thus \(\mathsf H_{\rm cross}=0\) whenever either group is empty.

Conditional on the right group, the finite-class estimate applies with maximal direction norm \(\mathsf K_{\rm row}(\boldsymbol v_{\rm right})\). Choosing each right-group sign to agree with its projection gives
\[
\max_{u_{\rm dir}\in\mathcal U_{\rm right}(\boldsymbol v_{\rm right})}
|\langle\mathsf S_\rho(\boldsymbol v_{\rm left}),u_{\rm dir}\rangle|
=
\sum_{i'=1}^{n_{\rm right}}
|\langle\mathsf S_\rho(\boldsymbol v_{\rm left}),v_{i'}^{\rm right}\rangle|.
\]
Consequently,
\[
\begin{aligned}
\mathbb E_{\rm left}\,
\mathsf H_{\rm cross}(\boldsymbol v_{\rm left},\boldsymbol v_{\rm right})
\leq{}&
n_{\rm left}\mathsf K_{\rm row}(\boldsymbol v_{\rm right})\\
&+4\mathbb E_{{\rm left},\rho}
\sum_{i'=1}^{n_{\rm right}}
|\langle\mathsf S_\rho(\boldsymbol v_{\rm left}),v_{i'}^{\rm right}\rangle|.
\end{aligned}
\]
% lean: CausalSmith.Experimentation.BivariateSobolevDesignCapacity.fixed_group_cross_integral_le_exact_rad

For any local sample, averaging over the fair signs cancels the off-diagonal terms:
\[
\mathbb E_\rho\|\mathsf S_\rho(\boldsymbol v_{\rm loc})\|_2^2
=\sum_{i=1}^{n_{\rm loc}}\|v_i^{\rm loc}\|_2^2.
\]
Thus
\[
\mathbb E\|\mathsf S_\rho(\boldsymbol v_{\rm loc})\|_2^2
=n_{\rm loc}r_0,
\qquad
\mathbb E\|\mathsf S_\rho(\boldsymbol v_{\rm loc})\|_2
\leq\sqrt{n_{\rm loc}r_0}.
\]
For every fixed \(u_{\rm dir}\in\mathbb R^{r_0}\), the projection estimate gives
\[
\mathbb E_{\rm right}
\sum_{i'=1}^{n_{\rm right}}
|\langle u_{\rm dir},v_{i'}^{\rm right}\rangle|
\leq n_{\rm right}\|u_{\rm dir}\|_2.
\]
Independence allows this estimate to be applied conditional on the left sample and its signs. Interchanging the integrable expectations yields
\[
\mathbb E
\sum_{i'=1}^{n_{\rm right}}
|\langle\mathsf S_\rho(\boldsymbol v_{\rm left}),v_{i'}^{\rm right}\rangle|
\leq
n_{\rm right}\sqrt{n_{\rm left}r_0}.
\]
After integrating the conditional cross-group bound, we obtain
\[
\mathbb E\mathsf H_{\rm cross}(\boldsymbol v_{\rm left},\boldsymbol v_{\rm right})
\leq
n_{\rm left}\mathbb E\mathsf K_{\rm row}(\boldsymbol v_{\rm right})
+4n_{\rm right}\sqrt{n_{\rm left}r_0}.
\]
% lean: CausalSmith.Experimentation.BivariateSobolevDesignCapacity.independent_group_cross_integral_le

\item We next bound the expected maximum signed norm.
Define the Euclidean unit ball
\[
\mathcal B_{\rm dir}
:=\{u_{\rm dir}\in\mathbb R^{r_0}:\|u_{\rm dir}\|_2\leq1\}.
\]
Norm duality and choosing the signs of the projections imply
\[
\mathsf K_{\rm row}(\boldsymbol v_{\rm loc})
=
\sup_{u_{\rm dir}\in\mathcal B_{\rm dir}}
\sum_{i=1}^{n_{\rm loc}}|\langle u_{\rm dir},v_i^{\rm loc}\rangle|.
\]
Indeed, for a fixed direction, the maximum absolute projection of a signed sum is the sum of the absolute projections; taking the supremum over unit directions gives its Euclidean norm.
% lean: CausalSmith.Experimentation.BivariateSobolevDesignCapacity.signingNormMax_eq_unit_projection_sum

For every nonempty finite \(\mathcal U_{\rm fin}\subset\mathcal B_{\rm dir}\),
\[
\max_{u_{\rm dir}\in\mathcal U_{\rm fin}}
|\langle\mathsf S_\rho(\boldsymbol v_{\rm loc}),u_{\rm dir}\rangle|
\leq\|\mathsf S_\rho(\boldsymbol v_{\rm loc})\|_2.
\]
The finite-class bound and the signed-norm first moment therefore give
\[
\mathbb E\max_{u_{\rm dir}\in\mathcal U_{\rm fin}}
\sum_{i=1}^{n_{\rm loc}}|\langle u_{\rm dir},v_i^{\rm loc}\rangle|
\leq n_{\rm loc}+4\sqrt{n_{\rm loc}r_0}.
\]
% lean: CausalSmith.Experimentation.BivariateSobolevDesignCapacity.finite_unit_projection_sum_integral_le

To pass to the whole unit ball, put
\[
\mathsf A_{\rm norm}
:=\mathbb E\sum_{i=1}^{n_{\rm loc}}\|v_i^{\rm loc}\|_2,
\qquad
0\leq\mathsf A_{\rm norm}<\infty.
\]
For an arbitrary \(\delta_{\rm net}>0\), set
\[
\varepsilon_{\rm net}
:=\frac{\delta_{\rm net}}{\mathsf A_{\rm norm}+1}>0.
\]
Compactness provides a nonempty finite set
\[
\mathcal U_{\rm net}\subset\mathcal B_{\rm dir},
\qquad
\forall u_{\rm dir}\in\mathcal B_{\rm dir}\ 
\exists u_{\rm net}\in\mathcal U_{\rm net}:
\|u_{\rm dir}-u_{\rm net}\|_2\leq\varepsilon_{\rm net}.
\]
For such a pair of directions,
\[
|\langle u_{\rm dir},v_i^{\rm loc}\rangle|
\leq
|\langle u_{\rm net},v_i^{\rm loc}\rangle|
+\varepsilon_{\rm net}\|v_i^{\rm loc}\|_2.
\]
Summing and taking the supremum gives
\[
\mathsf K_{\rm row}(\boldsymbol v_{\rm loc})
\leq
\max_{u_{\rm net}\in\mathcal U_{\rm net}}
\sum_{i=1}^{n_{\rm loc}}|\langle u_{\rm net},v_i^{\rm loc}\rangle|
+\varepsilon_{\rm net}\sum_{i=1}^{n_{\rm loc}}\|v_i^{\rm loc}\|_2.
\]
Hence
\[
\begin{aligned}
\mathbb E\mathsf K_{\rm row}(\boldsymbol v_{\rm loc})
&\leq n_{\rm loc}+4\sqrt{n_{\rm loc}r_0}
+\varepsilon_{\rm net}\mathsf A_{\rm norm}\\
&=n_{\rm loc}+4\sqrt{n_{\rm loc}r_0}
+\delta_{\rm net}\frac{\mathsf A_{\rm norm}}{\mathsf A_{\rm norm}+1}\\
&\leq n_{\rm loc}+4\sqrt{n_{\rm loc}r_0}+\delta_{\rm net}.
\end{aligned}
\]
Letting \(\delta_{\rm net}\) tend to zero proves
\[
\mathbb E\mathsf K_{\rm row}(\boldsymbol v_{\rm loc})
\leq n_{\rm loc}+4\sqrt{n_{\rm loc}r_0}.
\]
This also covers the empty sample.
% lean: CausalSmith.Experimentation.BivariateSobolevDesignCapacity.isotropic_signingNormMax_integral_le

Substitution into the preceding cross-group estimate now yields
\[
\begin{aligned}
\mathbb E\mathsf H_{\rm cross}(\boldsymbol v_{\rm left},\boldsymbol v_{\rm right})
\leq{}&
n_{\rm left}n_{\rm right}
+4n_{\rm left}\sqrt{n_{\rm right}r_0}\\
&+4n_{\rm right}\sqrt{n_{\rm left}r_0}.
\end{aligned}
\]
% lean: CausalSmith.Experimentation.BivariateSobolevDesignCapacity.independent_group_cross_integral_le_explicit

\item Introduce fair selectors, independent of the rows,
\[
\eta=(\eta_1,\ldots,\eta_n)\in\{0,1\}^n,
\qquad
\mathbb P(\eta_i=1)=\mathbb P(\eta_i=0)=\tfrac12,
\]
with independent coordinates. Define
\[
\begin{aligned}
\mathcal I_\eta&:=\{i\in\{1,\ldots,n\}:\eta_i=1\},&
\mathcal J_\eta&:=\{i\in\{1,\ldots,n\}:\eta_i=0\},\\
\mathsf N_{\rm sel}(\eta)&:=|\mathcal I_\eta|,&
\mathsf N_{\rm comp}(\eta)&:=|\mathcal J_\eta|,\\
\boldsymbol v_{\rm sel}&:=(v_i)_{i\in\mathcal I_\eta},&
\boldsymbol v_{\rm comp}&:=(v_i)_{i\in\mathcal J_\eta},
\end{aligned}
\]
where the subsequences are listed in increasing index order. For \(z\in\{-1,1\}^n\), define
\[
\begin{aligned}
\mathsf C_{\rm cross}(\boldsymbol v,z,\eta)
&:=\sum_{i\in\mathcal I_\eta}\sum_{i'\in\mathcal J_\eta}
z_i z_{i'}\langle v_i,v_{i'}\rangle\\
&=\left\langle
\sum_{i\in\mathcal I_\eta}z_i v_i,\,
\sum_{i'\in\mathcal J_\eta}z_{i'}v_{i'}
\right\rangle,\\
\mathsf M_{\rm cross}(\boldsymbol v,\eta)
&:=\max_{z\in\{-1,1\}^n}
|\mathsf C_{\rm cross}(\boldsymbol v,z,\eta)|.
\end{aligned}
\]
For fixed rows and selectors, the triangle inequality gives
\[
\begin{aligned}
|\mathsf C_{\rm cross}(\boldsymbol v,z,\eta)|
&\leq
\sum_{i\in\mathcal I_\eta}
\left|\left\langle
v_i,\sum_{i'\in\mathcal J_\eta}z_{i'}v_{i'}
\right\rangle\right|\\
&\leq\mathsf H_{\rm cross}(\boldsymbol v_{\rm sel},\boldsymbol v_{\rm comp}).
\end{aligned}
\]
For each fixed \(\eta\), the two subsequences are independent samples from the row law. The cross-group estimate thus implies
\[
\begin{aligned}
\mathbb E_{\boldsymbol v}\mathsf M_{\rm cross}(\boldsymbol v,\eta)
\leq{}&
\mathsf N_{\rm sel}(\eta)\mathsf N_{\rm comp}(\eta)\\
&+4\mathsf N_{\rm sel}(\eta)
\sqrt{\mathsf N_{\rm comp}(\eta)r_0}\\
&+4\mathsf N_{\rm comp}(\eta)
\sqrt{\mathsf N_{\rm sel}(\eta)r_0}.
\end{aligned}
\]
Since
\[
\mathsf N_{\rm sel}(\eta)+\mathsf N_{\rm comp}(\eta)=n,
\]
the square of their difference gives
\[
\mathsf N_{\rm sel}(\eta)\mathsf N_{\rm comp}(\eta)
\leq\frac{n^2}{4}.
\]
Both group sizes are at most \(n\), so
\[
\begin{aligned}
&4\mathsf N_{\rm sel}(\eta)
\sqrt{\mathsf N_{\rm comp}(\eta)r_0}
+4\mathsf N_{\rm comp}(\eta)
\sqrt{\mathsf N_{\rm sel}(\eta)r_0}\\
&\hspace{2em}\leq
4\bigl(\mathsf N_{\rm sel}(\eta)+\mathsf N_{\rm comp}(\eta)\bigr)
\sqrt{nr_0}
=4n\sqrt{nr_0}.
\end{aligned}
\]
Therefore every selector satisfies
\[
\mathbb E_{\boldsymbol v}\mathsf M_{\rm cross}(\boldsymbol v,\eta)
\leq\frac{n^2}{4}+4n\sqrt{nr_0}.
\]
% lean: CausalSmith.Experimentation.BivariateSobolevDesignCapacity.signingPartitionCrossMax_integral_le

Averaging this bound over the finite selector space and interchanging the integrable expectations gives
\[
\begin{aligned}
4\mathbb E_{\boldsymbol v}\mathbb E_\eta
\mathsf M_{\rm cross}(\boldsymbol v,\eta)
&=4\mathbb E_\eta\mathbb E_{\boldsymbol v}
\mathsf M_{\rm cross}(\boldsymbol v,\eta)\\
&\leq n^2+16n\sqrt{nr_0}.
\end{aligned}
\]
% lean: CausalSmith.Experimentation.BivariateSobolevDesignCapacity.signingPartitionCross_average_le

\item For distinct indices, independence of the selectors gives
\[
\mathbb P(i\in\mathcal I_\eta,\ i'\in\mathcal J_\eta)=\frac14,
\]
whereas an index cannot belong to both groups. Consequently, for every deterministic row tuple and signing,
\[
\mathsf O_{\rm row}(\boldsymbol v,z)
=4\mathbb E_\eta\mathsf C_{\rm cross}(\boldsymbol v,z,\eta).
\]
Taking absolute values and then maximizing yields
\[
\begin{aligned}
\mathsf M_{\rm off}(\boldsymbol v)
&=4\max_z
\left|\mathbb E_\eta\mathsf C_{\rm cross}(\boldsymbol v,z,\eta)\right|\\
&\leq4\max_z
\mathbb E_\eta|\mathsf C_{\rm cross}(\boldsymbol v,z,\eta)|\\
&\leq4\mathbb E_\eta\mathsf M_{\rm cross}(\boldsymbol v,\eta).
\end{aligned}
\]
Integrating and applying the partition-average estimate gives
\[
\mathbb E\mathsf M_{\rm off}(\boldsymbol v)
\leq4\mathbb E_{\boldsymbol v}\mathbb E_\eta
\mathsf M_{\rm cross}(\boldsymbol v,\eta)
\leq n^2+16n\sqrt{nr_0}.
\]
% lean: CausalSmith.Experimentation.BivariateSobolevDesignCapacity.signingOffDiagonal_integral_le_partition_average

The diagonal-subtraction inequality and its expected diagonal energy now give
\[
\begin{aligned}
\mathbb E\mathsf M_{\rm sign}(\boldsymbol v)
&\geq nr_0-\mathbb E\mathsf M_{\rm off}(\boldsymbol v)\\
&\geq nr_0-n^2-16n\sqrt{nr_0}.
\end{aligned}
\]
Together with the small-sample case, this proves the first assertion for every \(n\).
% lean: CausalSmith.Experimentation.BivariateSobolevDesignCapacity.isotropic_row_signing

\item Finally, assume \(r_0\geq4096n\). Multiplying this inequality by the nonnegative number \(r_0/4096\) gives
\[
nr_0\leq\frac{r_0^2}{4096}
=\left(\frac{r_0}{64}\right)^2,
\qquad
\sqrt{nr_0}\leq\frac{r_0}{64}.
\]
Multiplying by nonnegative factors also gives
\[
16n\sqrt{nr_0}\leq\frac{nr_0}{4},
\qquad
n^2\leq\frac{nr_0}{4096}.
\]
Hence
\[
n^2+16n\sqrt{nr_0}
\leq\left(\frac1{4096}+\frac14\right)nr_0
\leq\frac{nr_0}{2}.
\]
These inequalities include \(n=0\). Subtracting this error bound from the first assertion proves
\[
\mathbb E\mathsf M_{\rm sign}(\boldsymbol v)\geq\frac{nr_0}{2}.
\]
% lean: CausalSmith.Experimentation.BivariateSobolevDesignCapacity.isotropic_signing_error_le_half
\end{enumerate}
\end{proof}
